\subsection*{6.3 Triple Integrals}

\begin{conceptbox}
    The triple integral of a scalar function over a volume $V$ is defined as:
    \begin{equation*}
        \iiint_V f(x,y,z)\, dV
    \end{equation*}
\end{conceptbox}

\textbf{Expanded Explanation:}

A triple integral represents the accumulation of a quantity throughout a three-dimensional region.

To understand this, imagine dividing the volume $V$ into many tiny rectangular boxes of volume $\Delta V$. At each small box, the function $f(x,y,z)$ gives a value (such as density, temperature, or intensity).

We approximate the total accumulation by:
\begin{equation*}
    \sum f(x_i,y_i,z_i)\,\Delta V
\end{equation*}

As the size of these boxes shrinks to zero, the sum becomes the triple integral:
\begin{equation*}
    \iiint_V f(x,y,z)\, dV
\end{equation*}

\textbf{Interpretations in Engineering:}
\begin{itemize}
    \item If $f(x,y,z)=1$, the integral gives \textbf{volume}
    \item If $f(x,y,z)$ is density, the integral gives \textbf{mass}
    \item If $f(x,y,z)$ is energy density, the integral gives \textbf{total energy}
\end{itemize}

\textbf{Why Change Coordinates?}

For symmetric regions like spheres or cylinders, Cartesian coordinates become complicated. In such cases, we switch to:
\begin{itemize}
    \item Cylindrical coordinates
    \item Spherical coordinates
\end{itemize}

because they simplify both the region and the integrand.

\begin{examplebox}
    Evaluate $\iiint_V (x^2+y^2+z^2)\, dV$ over the sphere $x^2+y^2+z^2\le1$.

    \textbf{Step 1: Recognize symmetry}

    The region is a sphere centered at the origin:
    \[
        x^2+y^2+z^2 \le 1
    \]

    This suggests using \textbf{spherical coordinates}, where:
    \begin{align*}
        x & = r\sin\theta\cos\phi \\
        y & = r\sin\theta\sin\phi \\
        z & = r\cos\theta
    \end{align*}

    Most importantly:
    \[
        x^2+y^2+z^2 = r^2
    \]

    So the integrand simplifies significantly.

    \textbf{Step 2: Determine bounds}

    For a unit sphere:
    \begin{align*}
        0 \le r \le 1        \\
        0 \le \theta \le \pi \\
        0 \le \phi \le 2\pi
    \end{align*}

    \textbf{Step 3: Volume element}

    In spherical coordinates:
    \begin{equation*}
        dV = r^2\sin\theta\,dr\,d\theta\,d\phi
    \end{equation*}

    This factor $r^2\sin\theta$ accounts for the distortion when converting from Cartesian coordinates.

    \textbf{Step 4: Substitute into the integral}

    The integrand becomes:
    \[
        x^2+y^2+z^2 = r^2
    \]

    Thus:
    \begin{align*}
        \iiint_V (x^2+y^2+z^2)\, dV
         & = \int_0^{2\pi}\int_0^\pi\int_0^1 r^2 \cdot r^2\sin\theta\,dr\,d\theta\,d\phi \\
         & = \int_0^{2\pi}\int_0^\pi\int_0^1 r^4\sin\theta\,dr\,d\theta\,d\phi
    \end{align*}

    \textbf{Step 5: Separate the integrals}

    Because the variables are independent:
    \begin{align*}
        = \left(\int_0^1 r^4\,dr\right)
        \left(\int_0^\pi \sin\theta\,d\theta\right)
        \left(\int_0^{2\pi} d\phi\right)
    \end{align*}

    \textbf{Step 6: Evaluate each integral}

    \textbf{Radial integral:}
    \begin{align*}
        \int_0^1 r^4\,dr = \left[\frac{r^5}{5}\right]_0^1 = \frac{1}{5}
    \end{align*}

    \textbf{Angular ($\theta$) integral:}
    \begin{align*}
        \int_0^\pi \sin\theta\,d\theta = [-\cos\theta]_0^\pi = (-(-1)) - (-1) = 2
    \end{align*}

    \textbf{Azimuthal ($\phi$) integral:}
    \begin{align*}
        \int_0^{2\pi} d\phi = 2\pi
    \end{align*}

    \textbf{Step 7: Multiply results}

    \begin{align*}
        \frac{1}{5} \cdot 2 \cdot 2\pi = \frac{4\pi}{5}
    \end{align*}

    \textbf{Final Answer:}
    \[
        \boxed{\frac{4\pi}{5}}
    \]
\end{examplebox}